\section{Zadanie 2}
\autor{Hubert Jastrzębski}
Rzucamy kostką, aż nie zanotujemy dwóch szóstek pod rząd. Jaka jest oczekiwana liczba rzutów?

\textit{Wskazówka:} Odpowiedzią nie jest 36.

\subsection*{Rozwiązanie}
Zasada chyba (bo jeszcze nikt go nie spisał) podobna do tej z poprzedniego zadania,
ale nawet trochę prostsza.

Generalnie będziemy chcieli stworzyć układ równań, który pozwoli nam wyliczyć
oczekiwaną liczbę rzutów -- rozpisałem się trochę (+ w vs code ciężko czasem nie kliknąć tego tab), ale samo zadanie nie jest skomplikowane.

Zdefiniujmy sobie trzy stany:
\begin{itemize}
    \item Stan $N$ (neutralny) -- stan w którym w ostatnim rzucie nie wypadła szóstka (albo jest to pierwszy rzut).
    \item Stan $S$ (jedna szóstka) -- stan w którym w ostatnim rzucie wypadła szóstka, ale nie mamy jeszcze dwóch szóstek pod rząd.
    \item Stan $K$ (dwie szóstki, koniec) -- stan w którym mamy już dwie szóstki pod rząd.
\end{itemize}

Niech $X$ będzie zmienną losową, oznaczającą liczbę rzutów potrzebnych do zakończenia gry.

Niech:
\begin{itemize}
    \item $E_N = E(X | \text{ jesteśmy w stanie } N)$ -- oczekiwana liczba rzutów do zakończenia gry, jeśli zaczynamy w stanie $N$,
    \item $E_S = E(X | \text{ jesteśmy w stanie } S)$ -- oczekiwana liczba rzutów do zakończenia gry, jeśli zaczynamy w stanie $S$,
    \item $E_K = E(X | \text{ jesteśmy w stanie } K)$ -- oczekiwana liczba rzutów do zakończenia gry, jeśli zaczynamy w stanie $K$.
\end{itemize}

Oczywiście $E_K = 0$, bo gra się już skończyła.

Rozważmy teraz stan $N$. W tym stanie rzucamy kostką:
\begin{itemize}
    \item Z prawdopodobieństwem $\frac{1}{6}$ wypada szóstka i przechodzimy do stanu $S$.
    \item Z prawdopodobieństwem $\frac{5}{6}$ nie wypada szóstka i pozostajemy w stanie $N$.
\end{itemize}
W każdym przypadku wykonaliśmy jeden rzut, więc możemy zapisać równanie:
\[E_N = \frac{1}{6}(E_S + 1) + \frac{5}{6}(E_N + 1)\]
Podobnie jest ze stanem $S$:
\begin{itemize}
    \item Z prawdopodobieństwem $\frac{1}{6}$ wypada szóstka i przechodzimy do stanu $K$.
    \item Z prawdopodobieństwem $\frac{5}{6}$ nie wypada szóstka i wracamy do stanu $N$.
\end{itemize}
Znowu wykonaliśmy jeden rzut, więc mamy:
\[E_S = \frac{1}{6}(E_K + 1) + \frac{5}{6}(E_N + 1)\]

Teraz pozostaje tylko rozwiązać układ równań (bardzo ciężki):
\begin{align*}
E_N &= \frac{1}{6}(E_S + 1) + \frac{5}{6}(E_N + 1) \\
&= \frac{1}{6}(\frac{1}{6}(E_K + 1) + \frac{5}{6}(E_N + 1) + 1) + \frac{5}{6}(E_N + 1) \\
&= \frac{1}{6}(\frac{1}{6} + \frac{5}{6}E_N + \frac{5}{6} + 1) + \frac{5}{6}(E_N + 1) \\
&= \frac{2}{6} + \frac{5}{36}E_N + \frac{5}{6}E_N + \frac{5}{6} \\
&= \frac{7}{6} + \frac{35}{36}E_N \\
E_N &= \frac{7}{6} \cdot \frac{36}{1} = 42
\end{align*}

Jako iż $N$ to jest stan z którego rozpoczynamy grę, to odpowiedzią jest właśnie $E_N = 42$.

\subsection*{Odpowiedź}
\[
\boxed{\text{42}}
\]
